\documentclass[12pt]{article}
\usepackage[T1]{fontenc}
\usepackage[utf8]{inputenc}
\usepackage{lmodern}
\usepackage{textcomp}
\usepackage{enumitem}
\usepackage{geometry}
\geometry{margin=1in}
\begin{document}
\begin{center}
Answer Key - Biology - Graduate Level Assessment 10 \
\small (Answers are ordered exactly as in the question paper.)
\end{center}

\section*{Multiple Choice}
\begin{enumerate}[label=\arabic*.]
\item \textbf{Answer:} A
\\ \textit{Options:} A) High detrital biomass; B) High number of predatory species; C) High presence of fully matured plant species; D) High frequency of R-selected species; E) High rates of soil nutrient depletion
\\ \textit{Note:} Early-successional plant communities are characterized by rapid growth and high rates of primary productivity, leading to a significant accumulation of organic matter in the form of detritus, which contributes to high detrital biomass.
\item \textbf{Answer:} D
\\ \textit{Options:} A) Acquired characteristics are inherited through a blending of parental traits; B) Acquired characteristics are inherited but do not influence evolution; C) Acquired characteristics are inherited randomly, regardless of their use; D) Acquired characteristics cannot be inherited; E) Acquired characteristics are inherited based on use and disuse
\\ \textit{Note:} Lamarck's theory posits that traits acquired during an organism's lifetime can be passed on to offspring, but modern genetics has demonstrated that only genetic traits can be inherited, thus refuting the notion of inheritance of acquired characteristics.
\item \textbf{Answer:} B
\\ \textit{Options:} A) Lack of available organs; B) Graft rejection due to immunological reaction; C) Incompatibility of blood types; D) Allergic reactions to anesthesia; E) Insufficient surgical expertise
\\ \textit{Note:} Graft rejection occurs because the recipient's immune system identifies the transplanted organ as foreign, leading to an immunological response that attacks the graft, making organ transplants generally unsuccessful without immunosuppressive therapy.
\item \textbf{Answer:} A
\\ \textit{Options:} A) The body controls water excretion through sweat glands.; B) Water excretion is determined by the concentration of electrolytes in the body's cells.; C) The kidneys control water excretion by changing the size of the glomerulus.; D) Water excreted is controlled by the hormone insulin.; E) The body adjusts water excretion by altering the pH level of the blood.
\\ \textit{Note:} The body controls water excretion through sweat glands as they play a crucial role in thermoregulation and fluid balance by allowing excess water to be lost through perspiration, thereby helping to maintain homeostasis.
\item \textbf{Answer:} A
\\ \textit{Options:} A) Telophase; B) G\_2; C) Prophase; D) Interphase; E) Metaphase
\\ \textit{Note:} The correct answer is telophase because, during this phase, the chromosomes that were duplicated during the S phase of interphase are segregated and begin to de-condense, resulting in the cell preparing to divide with a doubled quantity of DNA.
\end{enumerate}

\section*{True / False}
\begin{enumerate}[label=\arabic*.]
\setcounter{enumi}{5}
\item \textbf{Answer:} True
\\ \textit{Options:} A) True; B) False
\\ \textit{Note:} In manikin laypersons could insert LMAS in the correct direction after onsite instruction by a simple manual with a high success rate. This indicates some basic procedural understanding and intellectual transfer in principle. Operating errors (n = 91) were frequently not recognized and corrected (n = 77). Improvements in labeling and the quality of instructional photographs may reduce individual error and may optimize understanding.
\item \textbf{Answer:} True
\\ \textit{Options:} A) True; B) False
\\ \textit{Note:} LA for adrenal masses larger than 7 cm is a safe and feasible technique, offering successful outcome in terms of intraoperative and postoperative morbidity, hospital stay and cosmesis for patients; it seems to replicate open surgical oncological principles demonstrating similar outcomes as survival rate and recurrence rate, when adrenal cortical carcinoma were treated. The main contraindication for this approach is the evidence, radiologically and intraoperatively, of local infiltration of periadrenal tissue.
\item \textbf{Answer:} False
\\ \textit{Options:} A) True; B) False
\\ \textit{Note:} Overall, there was a poor knowledge of the side effects of ACE-I. This may account for the increased referrals for chronic cough and angioedema.
\item \textbf{Answer:} True
\\ \textit{Options:} A) True; B) False
\\ \textit{Note:} The results confirm that c-kit is vastly expressed in uveal melanoma, suggest that the c-kit molecular pathway may be important in uveal melanoma growth, and point to its use as a target for therapy with STI 571.
\item \textbf{Answer:} True
\\ \textit{Options:} A) True; B) False
\\ \textit{Note:} Our findings reveal a significant decrease in ADMA levels of ex-ELBW subjects compared to C, underlining a probable correlation with preterm birth and low birth weight. Taken together, these results may underlie the onset of early circulatory dysfunction predictive of increased cardiovascular risk.
\end{enumerate}

\section*{Long Form Response}
\begin{enumerate}[label=\arabic*.]
\setcounter{enumi}{10}
\item \textbf{Answer:} Blood pressure readings taken from the upper arm are typically higher than those taken from the leg due to the physiological effects of distance from the heart. As blood travels through the circulatory system, it encounters resistance from blood vessel walls, leading to a gradual decrease in pressure. This reduction is influenced by factors such as vessel diameter and the cumulative effect of resistance encountered in the arterial tree. Consequently, the upper arm, being closer to the heart, experiences higher pressure compared to the leg, where blood has to traverse a longer distance and overcome more vascular resistance. Thus, the differences in blood pressure readings can be attributed primarily to the distance from the heart and the associated hemodynamic changes.
\\ \textit{Note:} correct because it accurately describes how blood pressure decreases with distance from the heart due to increased resistance in the vascular system, leading to higher readings in the upper arm compared to the leg.
\item \textbf{Answer:} Before the works of Charles Lyell and Charles Darwin, the prevailing scientific beliefs posited that the Earth was only a few thousand years old, largely based on interpretations of biblical chronology. Additionally, the concept of species was rooted in the idea of fixity, meaning that populations were viewed as unchanging entities created in their current forms by a divine being. These views directly influenced Lyell and Darwin's theories of evolution by necessitating a paradigm shift; they provided the foundation for questioning the static nature of species and the geological timescale. Lyell's uniformitarianism emphasized gradual geological processes over immense timeframes, which allowed for the possibility of biological change, while Darwin's theory of natural selection built upon the idea that species could evolve over time in response to environmental pressures. Thus, the earlier beliefs about Earth's young age and species' constancy set the stage for the revolutionary ideas that would emerge in the field of evolutionary biology.
\\ \textit{Note:} correct because it accurately describes the historical scientific beliefs regarding the Earth's young age and the fixity of species, illustrating how these views created a foundational context that Lyell and Darwin needed to challenge in order to propose their revolutionary theories of geological time and evolution.
\end{enumerate}

\vfill
\noindent\textit{End of Answer Key}
\end{document}